\section{Related Work and Novelty Boundary}
\label{sec:related}

\paragraph{Regression testing and change impact.}
Regression-test selection has long been studied as a way to reuse tests after
software modification \cite{RothermelHarrold1996,YooHarman2012}.  Code-based
change-impact analysis identifies which behavior may be affected by a change
\cite{LiSunLeungZhang2013}.  Our model treats sound impact information as side
information in a statistical query problem rather than proposing a new static
impact analysis.

\paragraph{Fairness testing and auditing.}
Themis established automated software fairness testing as a software-engineering
problem \cite{GalhotraBrunMeliou2017}.  Active fairness auditing gives theoretical
query-efficient algorithms for assessing fixed learned systems
\cite{YanZhang2022}.  We do not claim novelty for fairness testing or active
fairness auditing in isolation; our target is version-aware stochastic regression.

\paragraph{Property testing.}
Property testing provides the surrounding query-complexity framework.
Optimal monotonicity/Lipschitz testing on hypercubes and hypergrids is a canonical
example of graph-structured local checks determining a global property
\cite{ChakrabartySeshadhri2013}; tolerant testing can be much harder than ordinary
testing \cite{FischerFortnow2006}; and adaptivity hierarchies are known
\cite{CanonneGur2018}.  Our property and access model are different: observations
are stochastic executions of two evolving software versions, and old-version and
impact information are first-class inputs.

\paragraph{Adaptive information lower bounds.}
Adaptive fixed-confidence sampling is commonly analyzed through transcript-level
change-of-measure arguments.  We use the generic information-theoretic technique
of Kaufmann, Capp{\'e}, and Garivier \cite{KaufmannCappeGarivier2016} and
specialize the alternatives to counterfactual patch effects on graph edges.

\paragraph{Pairwise experimental design.}
Optimal design for selected pairwise mean comparisons has a substantial
statistical literature, including recent graph-based formulations
\cite{SinghSinghDavidov2025}.  Therefore the existence of a graph allocation
problem by itself is not a novelty claim.  The role of $\Psi$ here is as the core
information quantity arising from a regression-specific coupled execution oracle;
the intended contribution is the complete software-evolution query model and its
upper/lower complexity theory.

\paragraph{Incremental and probabilistic verification.}
Incremental verification seeks to reuse proof information from earlier program
versions \cite{YinEtAl2025}.  Couplings are a standard tool for relational
reasoning about probabilistic programs \cite{BartheEtAl2017}.  We combine the
version-reuse viewpoint with statistical testing of a fairness property rather
than attempting full formal equivalence verification.

\paragraph{Nondeterministic software testing.}
Flaky tests demonstrate that stochastic and environment-dependent execution is a
practical testing concern.  Recent OOPSLA work controls nondeterministic APIs to
expose flakiness \cite{YuanLinShi2026}.  Our theory abstracts such uncertainty as
an execution oracle and asks how it changes the fundamental cost of fairness
regression certification.
